%------------------------------------------------------------------------------%
\documentclass[fleqn,utf8,aspectratio=169,12pt]{beamer}

\usepackage{integracao_e_series}
\usepackage{systeme}
\usepackage{polynom}

\title{Integração por Frações Parciais}
\subtitle{Caso 1 -- Fatores reais distintos}

%------------------------------------------------------------------------------%
\begin{document}

\addtitle

%------------------------------------------------------------------------------%
\section{Objetivo e Preliminares}

%------------------------------------------------------------------------------%
\begin{frame}{Objetivo}
  Integrar funções da forma
  \[
    f(x) = \frac{P(x)}{Q(x)}
  \]
  onde $P$ e $Q$ são polinômios
\end{frame}

%------------------------------------------------------------------------------%
\begin{frame}{Fatoração de Polinômios}
  Se o polinômio $P$, de grau $n$,
  \[
    P(x) = ax^n + bx^{n-1} + \cdots + rx + s
  \]
  \pause
  Possui $n$ raízes reais distintas $x_i$
  \[
    P(x) = a(x-x_1)(x-x_2)\cdots(x-x_n)
  \]
  \pause
  Se a raiz $x_k$ se repete $p$ vezes
  \[
    P(x) = a(x-x_1)(x-x_2)\cdots(x-x_k)^p
  \]
  \pause
  Se $P$ possui uma raiz complexa
  \[
    P(x) = a(x-x_1)(x-x_2)\cdots(x^2+bx+c)
  \]
\end{frame}

%------------------------------------------------------------------------------%
\begin{frame}{Integral Relevante}
  \begin{align*}
    \onslide<+->{F(x) & = \int \frac{m}{ax+b}\, dx}
    \onslide<+->{& u = ax+b \qquad du = adx \\}
    \onslide<+->{& = \int \frac{m}{u}\, \frac{du}{a}   \\}
    \onslide<+->{& = \frac{m}{a}\int \frac{1}{u}\, du  \\}
    \onslide<+->{& = \frac{m}{a}\ln\abs{u} + c \\}
    \onslide<+->{& = \frac{m}{a}\ln\abs{x+a} + c}
  \end{align*}
\end{frame}

%------------------------------------------------------------------------------%
%------------------------------------------------------------------------------%
\begin{question}
  Integrar \(\int \frac{x^3+x}{x-1}\,dx\)
\end{question}

%------------------------------------------------------------------------------%
\begin{resolution}{Solução}
  \[
    f(x) = \frac{P(x)}{Q(x)}
    \pause
    \qquad\qquad
    P(x) = x^3+x
    \pause
    \qquad\qquad
    Q(x) = x-1
  \]

  \pause
  Como o grau de $P$ é maior do que o de $Q$ fazemos a divisão
\end{resolution}

%------------------------------------------------------------------------------%
\begin{resolutiont}{Solução}
 \vspace{1.5\baselineskip}
\begin{columns}[t]
\column{0.5\linewidth}
  \only<+>{\polylongdiv[style=D, stage= 1]{x^3+x}{x-1}}%
  \only<+>{\polylongdiv[style=D, stage= 2]{x^3+x}{x-1}}%
  \only<+>{\polylongdiv[style=D, stage= 3]{x^3+x}{x-1}}%
  \only<+>{\polylongdiv[style=D, stage= 4]{x^3+x}{x-1}}%
  \only<+>{\polylongdiv[style=D, stage= 5]{x^3+x}{x-1}}%
  \only<+>{\polylongdiv[style=D, stage= 6]{x^3+x}{x-1}}%
  \only<+>{\polylongdiv[style=D, stage= 7]{x^3+x}{x-1}}%
  \only<+>{\polylongdiv[style=D, stage= 8]{x^3+x}{x-1}}%
  \only<+>{\polylongdiv[style=D, stage= 9]{x^3+x}{x-1}}%
  \only<+->{\polylongdiv[style=D, stage=10]{x^3+x}{x-1}}%
\column{0.5\linewidth}
  \begin{align*}
    \only<+->{P(x)}
    \only<+->{& = x^3 + x} \\
    \only<+->{& = \left(x^2 + x + 2\right)(x-1) + 2}
  \end{align*}
\end{columns}
\end{resolutiont}

%------------------------------------------------------------------------------%
\begin{resolution}{Solução}
  \begin{align*}
    \onslide<+->{f(x)}
    \onslide<+->{& = \frac{P(x)}{Q(x)}                                          }\\
    \onslide<+->{& = \frac{x^3+x}{x-1}                                          }\\
    \onslide<+->{& = \frac{\left(x^2 + x + 2\right)(x-1) + 2}{x-1}              }\\
    \onslide<+->{& = \frac{\left(x^2 + x + 2\right)(x-1)}{x-1} + \frac{2}{x-1}  }\\
    \onslide<+->{& = x^2 + x + 2 + \frac{2}{x-1}}
  \end{align*}
\end{resolution}

%------------------------------------------------------------------------------%
\begin{resolution}{Solução}
  \begin{align*}
    \onslide<+->{\int \frac{x^3+x}{x-1} \,dx}
    \onslide<+->{& = \int x^2 + x + 2 + \frac{2}{x-1} \, dx \\[2mm]}
    \onslide<+->{& = \int x^2 \, dx
                   + \int x   \, dx
                   + \int 2   \, dx
                   + \int \frac{2}{x-1} \, dx \\[2mm]}
    \onslide<+->{& = \frac{x^3}{3}
      + \frac{x^2}{2}
      + 2x
      + 2 \ln\abs{x-1}
      + c} \\[2mm]
    \onslide<+->{& \color{gray} = \frac{x^3}{3}
      + \frac{x^2}{2}
      + 2x
      + \ln(x-1)^2
      + c}
  \end{align*}
\end{resolution}

%------------------------------------------------------------------------------%


%------------------------------------------------------------------------------%
\section{Integração por Frações Parciais}

%------------------------------------------------------------------------------%
\begin{frame}{Integração por Frações Parciais}
  Integrar \(f(x)=\frac{P(x)}{Q(x)}\) onde $P$ e $Q$ são polinômios
  \pause
  \begin{enumerate}[<+->]
    \setlength\itemsep{1em}
    \item Se o grau de $P$ for maior ou igual que o grau de $Q$, dividir $P$ por $Q$

    \item Fatorar $Q(x)$

    \item Expandir $f$ em frações parciais
  \end{enumerate}
\end{frame}

%------------------------------------------------------------------------------%
\begin{frame}{Frações Parciais -- Caso 1}
  $Q(x)$ é o produto de $n$ fatores lineares distintos
  \[
    Q(x) =
    \left( a_1 x + b_1 \right)
    \left( a_2 x + b_2 \right) \cdots
    \left( a_n x + b_n \right)
  \]

  \pause
  Expansão em Frações Parciais
  \[
    \frac{P(x)}{Q(x)}
    = \frac{A_1}{ a_1 x + b_1 }
    + \frac{A_2}{ a_2 x + b_2 }
    + \cdots
    + \frac{A_m}{ a_n x + b_n }
  \]
\end{frame}

%------------------------------------------------------------------------------%
\section{Exemplos}

%------------------------------------------------------------------------------%
\begin{question}
  Calcule
  \(
    \int \frac{3x + 5}{(x+1)(x+2)} \, dx
  \)
\end{question}

%------------------------------------------------------------------------------%
\begin{resolution}{Solução}
  O grau do numerador
  \[
    P(x) = 3x + 5
  \]
  \pause
  é menor do que o do denominador
  \[
    Q(x) = (x+1)(x+2)
  \]
  \pause
  O denominador já está fatorado

  \pause
  Por frações parciais temos
  \[
    \pause
    f(x)
    \;=\; \frac{P(x)}{Q(x)}
    \pause
    \;=\; \frac{3x + 5}{(x+1)(x+2)}
    \pause
    \;=\; \frac{A}{x+1} + \frac{B}{x+2}
  \]
\end{resolution}

%------------------------------------------------------------------------------%
\begin{resolution}{Solução}
  \begin{align*}
    \onslide<+->{\frac{3x + 5}{(x+1)(x+2)} & =  \frac{A}{x+1} + \frac{B}{x+2}      \\[3mm]}
    \onslide<+->{3x + 5 & = \left(\frac{A}{x+1} + \frac{B}{x+2}\right)(x+1)(x+2)   \\[1mm]}
    \onslide<+->{& = A(x+2) + B(x+1)                                               \\[1mm]}
    \onslide<+->{& = Ax + 2A + Bx + B                                              \\[1mm]}
    \onslide<+->{& = (A+B)x + (2A + B)}
  \end{align*}
\end{resolution}

%------------------------------------------------------------------------------%
\begin{resolution}{Solução}
\begin{columns}
\column{0.5\linewidth}
  \onslide<+->{Para que os polinômios sejam iguais
  \[
    3x + 5 = (A+B)x + (2A + B)
  \]}
  \onslide<+->{os coeficientes precisam ser iguais}
  \[
    \onslide<+->{
    \systeme[sort={A,B}]{
      A + B = 3,
      2A + B = 5
    }}
  \]
\column{0.5\linewidth}
  \[
    \onslide<+->{B = 3 - A}
  \]
  \vspace{-1.5\baselineskip}
  \begin{align*}
    \onslide<+->{2A + B     & = 5     \\}
    \onslide<+->{2A + 3 - A & = 5     \\}
    \onslide<+->{A          & = 5 - 3}
    \onslide<+->{             = 2}
  \end{align*}
  \vspace{-1.5\baselineskip}
  \begin{align*}
    \onslide<+->{B & = 3 - A \\}
    \onslide<+->{  & = 3 - 2   }
    \onslide<+->{    = 1}
  \end{align*}
\end{columns}
\end{resolution}

%------------------------------------------------------------------------------%
\begin{resolution}{Solução}
\begin{align*}
  \onslide<+->{f(x)}
  \onslide<+->{& = \frac{3x + 5}{(x+1)(x+2)}           \\[2mm]}
  \onslide<+->{& = \frac{A}{x+1} + \frac{B}{x+2}       \\[2mm]}
  \onslide<+->{& = \frac{2}{x+1} + \frac{1}{x+2}}
\end{align*}
\end{resolution}

%------------------------------------------------------------------------------%
\begin{resolution}{Solução}
  \begin{align*}
    \onslide<+->{F(x)}
    \onslide<+->{& = \int \frac{3x + 5}{(x+1)(x+2)} dx              \\[2mm]}
    \onslide<+->{& = \int \frac{2}{x+1} + \frac{1}{x+2} dx          \\[2mm]}
    \onslide<+->{& = 2\int \frac{1}{x+1} dx + \int \frac{1}{x+2} dx \\[2mm]}
    \onslide<+->{& = 2\ln\abs{x+1} + \ln\abs{x+2} + c               \\[2mm]}
    \onslide<+->{& = \color{gray} \ln\left((x+1)^2\,\abs{x+2}\right) + c}
  \end{align*}
\end{resolution}

%------------------------------------------------------------------------------%

%------------------------------------------------------------------------------%
\begin{question}
  Calcule
  \(
    \int \frac{x-3}{x(x-1)(x+1)} \,dx
  \)
\end{question}

%------------------------------------------------------------------------------%
\begin{resolution}{Solução}
  O grau do numerador
  \[
    P(x) = x - 2
  \]
  \pause
  é menor do que o do denominador
  \[
    Q(x) = x(x-1)(x+1)
  \]
  \pause
  O denominador já está fatorado

  \pause
  Por frações parciais temos
  \[
    \pause
    f(x)
    \;=\; \frac{P(x)}{Q(x)}
    \pause
    \;=\; \frac{x-3}{x(x-1)(x+1)}
    \pause
    \;=\; \frac{A}{x} + \frac{B}{x-1} + \frac{C}{x+1}
  \]
\end{resolution}

%------------------------------------------------------------------------------%
\begin{resolution}{Solução}
  \begin{align*}
    \onslide<+->{\frac{x-3}{x(x-1)(x+1)}}
    \onslide<+->{& = \frac{A}{x}
                   + \frac{B}{x-1}
                   + \frac{C}{x+1}
    \\[1mm]}
    \onslide<+->{x-3
                 & = \left(
                     \frac{A}{x}
                   + \frac{B}{x-1}
                   + \frac{C}{x+1}
                   \right)
                   x(x-1)(x+1)
    \\[1mm]}
    \onslide<+->{& = A(x-1)(x+1)
                   + Bx(x+1)
                   + Cx(x-1)
    \\}
    \onslide<+->{& = A\left(x^2 + 1\right)
                   + B\left(x^2 + x\right)
                   + C\left(x^2 - x\right)
    \\}
    \onslide<+->{& = x^2 \left( A + B + C \right)
                   + x   \left(B-C\right)
                   - A}
  \end{align*}
\end{resolution}

%------------------------------------------------------------------------------%
\begin{resolution}{Solução}
  \[
    \onslide<+->{
    x - 3
    \;=\; x^2 \left(A + B + C \right)
    \;+\; x   \left(B - C\right)
    \;-\; A}
  \]

\vspace{-\baselineskip}
\begin{columns}[t]
\column{0.3\linewidth}
  \[
    \onslide<+->{\systeme[sort={A,B,C}]{
      A + B + C = 0,
          B - C = 1,
      A         = 3
    }}
  \]

  \smallskip
  \[
    \onslide<+->{B = 1 + C}
  \]
\column{0.4\linewidth}
  \begin{align*}
    \onslide<+->{A + B + C     & = 0  \\}
    \onslide<+->{3 + 1 + C + C & = 0  \\}
    \onslide<+->{2C            & = -4 \\}
    \onslide<+->{C             & = -2}
  \end{align*}
  \[
    \onslide<+->{B = 1 - 2 = -1}
  \]
\end{columns}
\end{resolution}

%------------------------------------------------------------------------------%
\begin{resolution}{Solução}
\begin{align*}
  \onslide<+->{\frac{x-3}{x(x-1)(x+1)}
  & = \frac{A}{x}
    + \frac{B}{x-1}
    + \frac{C}{x+1} \\[2mm]}
  \onslide<+->{& = \frac{3}{x}
    + \frac{-1}{x-1}
    + \frac{-2}{x+1} \\[2mm]}
  \onslide<+->{& = \frac{3}{x}
    - \frac{1}{x-1}
    - \frac{2}{x+1}}
\end{align*}
\end{resolution}

%------------------------------------------------------------------------------%
\begin{resolution}{Solução}
\begin{align*}
  \onslide<+->{\int \frac{x-3}{x(x-1)(x+1)} \,dx}
  \onslide<+->{& = \int \frac{3}{x}
    - \frac{1}{x-1}
    - \frac{2}{x+1} \,dx \\[2mm]}
  \onslide<+->{& = \int \frac{3}{x} \,dx
    - \int \frac{1}{x-1} \,dx
    - 2 \int \frac{1}{x+1} \,dx \\[2mm]}
  \onslide<+->{& = 3\ln\abs{x}
    - \ln\abs{x-1}
    - 2 \ln\abs{x+1}
    + k \\[2mm]}
  \onslide<+->{&\color{gray}
    = \ln\abs{\frac{x^3}{(x-1)(x+1)^2}}
    + k}
\end{align*}
\end{resolution}

%------------------------------------------------------------------------------%

%------------------------------------------------------------------------------%
\begin{question}
  Calcule
  \(
    \int \frac{x^2 + 2x - 1}{2x^3 + 3x^2 - 2x} \, dx
  \)
\end{question}

%------------------------------------------------------------------------------%
\begin{resolution}{Solução}
  O grau do numerador
  \[
    P(x) = x^2 + 2x - 1
  \]
  \pause
  é menor do que o do denominador
  \[
      Q(x) = 2x^3 + 3x^2 -2x
  \]
  \pause
  Precisamos fatorar o denominador
\end{resolution}

%------------------------------------------------------------------------------%
\begin{resolution}{Fatorando $Q$}
  \[
    Q(x)
    \;=\; 2x^3 + 3x^2 - 2x
    \pause
    \;=\; x\left(2x^2 + 3x - 2\right)
  \]
  \[
    \pause
    x_1 = 0
    \pause
    \qquad
    \text{ ou }
    \qquad
    2x^2 + 3x - 2 = 0
  \]
  \[
    \pause
    \Delta
    = b^2 -4ac
    \pause
    = 3^2 - 4\times 2 \times (-2)
    \pause
    = 25
  \]
  \[
    \pause
    x_2
    = \frac{-b + \sqrt{\Delta}}{2a}
    \pause
    = \frac{-3 + 5}{2\times 2}
    \pause
    = \frac{2}{4}
    \pause
    = \frac{1}{2}
  \]
  \[
    \pause
    x_3
    \pause
    = \frac{-b - \sqrt{\Delta}}{2a}
    \pause
    = \frac{-3 - 5}{2\times 2}
    \pause
    = \frac{-8}{4}
    \pause
    = -2
  \]
  \pause
  \[
    Q(x)
    \pause
    = x\, 2\left(x- \frac{1}{2}\right)\left(x+2\right)
    \pause
    = x\left(2x-1\right)\left(x+2\right)
  \]
\end{resolution}

%------------------------------------------------------------------------------%
\begin{resolution}{Por frações parciais temos}
  \[
    \onslide<+->{\frac{x^2+2x-1}{2x^3+3x^2-2x}}
    \onslide<+->{\;=\; \frac{x^2+2x-1}{x\left(2x-1\right)\left(x+2\right)}}
    \onslide<+->{\;=\; \frac{A}{x}
                   + \frac{B}{2x-1}
                   + \frac{C}{x+2}}
  \]
  \begin{align*}
   \onslide<+->{ x^2+2x-1
    & = \left(
        \frac{A}{x}
      + \frac{B}{2x-1}
      + \frac{C}{x+2}
      \right)
      x\left(2x-1\right)\left(x+2\right)}
    \\[1mm]
    \onslide<+->{& = A(2x-1)(x+2)+Bx(x+2)+ Cx(2x-1)}
    \\[1mm]
    \onslide<+->{& = A \left( 2x^2 + 3x -2 \right)
                   + B \left( x^2 + 2x     \right)
                   + C \left( 2x^2 - x     \right)}
    \\[1mm]
    \onslide<+->{& = (2A+B +2C)\,x^2 \,+\, (3A+2B-C)\,x \,-\, 2A}
  \end{align*}
\end{resolution}

%------------------------------------------------------------------------------%
\begin{resolution}{Solução}
  \onslide<+->{\[
    x^2+2x-1 \;=\; (2A+B +2C)x^2 \;+\; (3A+2B-C)x \;-\; 2A
  \]}

  \vspace{-2.5\baselineskip}
  \begin{columns}[t]
  \column{0.28\linewidth}
    ~
    \vspace{\baselineskip}

    \onslide<+->{\(
      \systeme[sort={A,B,C}]{
        2A + B +2C = 1,
        3A +2B - C = 2,
        2A         = 1
      }
    \)}

    \onslide<+->{\[
      \qquad
      A = \frac{1}{2}
    \]}
  \column{0.28\linewidth}
    \begin{align*}
      \onslide<+->{2A + B + 2C =  1   \\}
      \onslide<+->{1 +  B + 2C =  1   \\}
      \onslide<+->{B           = -2C    }
    \end{align*}
  \column{0.35\linewidth}
    \begin{align*}
      \onslide<+->{\frac{3}{2} + 2B -  C & = 2              \\}
      \onslide<+->{              2B -  C & =  \frac{1}{2}   \\}
      \onslide<+->{                  -5C & =  \frac{1}{2}   \\}
      \onslide<+->{                    C & = -\frac{1}{10}    }
    \end{align*}
  \end{columns}
  \vspace{-\baselineskip}
  \[
    \onslide<+->{A =  \frac{1}{2}  \qquad }
    \onslide<+->{B =  \frac{1}{5}  \qquad }
    \onslide<+->{C = -\frac{1}{10}        }
  \]
\end{resolution}

%------------------------------------------------------------------------------%
\begin{resolution}{Solução}
  \begin{align*}
  \onslide<+->{\int \frac{x^2+2x-1}{2x^3+3x^2-2x} dx}
  \onslide<+->{& = \int
      \frac{1}{2x}
    + \frac{1}{5(2x-1)}
    - \frac{1}{10(x+2)} \,dx}
  \\[2mm]
  \onslide<+->{& = \frac{1}{2}  \int \frac{1}{x} \,dx
    + \frac{1}{5}  \int \frac{1}{2x-1} \,dx
    - \frac{1}{10} \int \frac{1}{x+2} \,dx }
  \\[2mm]
  \onslide<+->{& = \frac{1}{2}  \ln\abs{x}
    + \frac{1}{10} \ln\abs{2x-1}
    - \frac{1}{10} \ln\abs{x+2}
    + c}
  \\[2mm]
  \onslide<+->{
    & \color{gray} = \ln\sqrt{\abs{x}}
    + \ln\sqrt[10]{\abs{2x-1}}
    - \ln\sqrt[10]{\abs{x+2}}
    + c}
  \\[1mm]
  \onslide<+->{
    & \color{gray} = \ln\left(\sqrt{\abs{x}}
    \sqrt[10]{\abs{\frac{2x-1}{x+2}}}\right)
    + c}
\end{align*}
\end{resolution}

%------------------------------------------------------------------------------%

%------------------------------------------------------------------------------%
\begin{question}
  Calcule
  \(
    \int \frac{\sin(x)}{\cos^2(x)-3\cos(x)}dx
  \)
\end{question}

%------------------------------------------------------------------------------%
\begin{resolution}{Solução}
  \[
    F(x) = \int \frac{\sin(x)}{\cos^2(x)-3\cos(x)}dx
  \]

  \[
    \pause
    u  =  \cos(x)
    \qquad
    \pause
    du = -\sin(x) dx
    \pause
    \qquad
    \sin(x)dx = -du
  \]

  \pause
  \[
    F(x)
    \;=\; \int \frac{1}{u^2-3u}(-du)
    \pause
    \;=\; -\int \frac{1}{u(u-3)}du
  \]

  \pause
  \[
    \frac{1}{u(u-3)}
    = \frac{A}{u} + \frac{B}{u-3}
  \]
\end{question}

%------------------------------------------------------------------------------%
\begin{resolution}{Solução}
\begin{columns}
\column{0.55\linewidth}
  \begin{align*}
    \onslide<+->{\frac{1}{u(u-3)} }
    \onslide<+->{  & = \frac{A}{u} + \frac{B}{u-3}        \\[3mm]}
    \onslide<+->{1 & = \left(\frac{A}{u} + \frac{B}{u-3}\right)\, u(u-3) \\}
    \onslide<+->{  & = A(u - 3) + Bu                                     \\}
    \onslide<+->{  & = (A + B)u - 3A                                       }
  \end{align*}
\column{0.45\linewidth}
  \[
    \onslide<+->{\systeme[sort={A,B}]{
      A + B = 0,
      -3A = 1
    }}
  \]

  \[
    \onslide<+->{A = -\frac{1}{3}}
  \]

  \[
    \onslide<+->{B = -A =  \frac{1}{3}}
  \]
\end{columns}
\end{resolution}

%------------------------------------------------------------------------------%
\begin{resolution}{Solução}
\begin{align*}
  \onslide<+->{\int \frac{1}{u(u-3)}du
  & = \int \frac{-\sfrac{1}{3}}{u} + \frac{\sfrac{1}{3}}{u-3} \,du \\[2mm]}
  \onslide<+->{& = -\frac{1}{3} \int \frac{1}{u}   \,du
      +\frac{1}{3} \int \frac{1}{u-3} \,du \\[2mm]}
  \onslide<+->{& = - \frac{1}{3} \ln\abs{u}
      + \frac{1}{3} \ln\abs{u-3} + c \\[2mm]}
  \onslide<+->{& = \frac{1}{3} \ln \abs{\frac{u-3}{u}} + c }
\end{align*}
\end{resolution}

%------------------------------------------------------------------------------%
\begin{resolution}{Solução}
  \begin{align*}
    \onslide<+->{\int \frac{\sin(x)}{\cos^2(x)-3\cos(x)}dx
    & = -\int \frac{1}{u(u-3)}du \\[2mm]}
    \onslide<+->{& = -\frac{1}{3} \ln \abs{\frac{u-3}{u}} + c \\[2mm]}
    \onslide<+->{& =  \frac{1}{3} \ln \abs{\frac{u}{u-3}} + c \\[2mm]}
    \onslide<+->{& =  \frac{1}{3} \ln \abs{\frac{\cos(x)}{\cos(x)-3}} + c}
  \end{align*}
\end{resolution}

%------------------------------------------------------------------------------%


%------------------------------------------------------------------------------%
\section{Lista Mínima}

%------------------------------------------------------------------------------%
\begin{frame}{Lista Mínima}
  Estudar a Seção 4.5 da Apostila

  Exercícios:

  \vfill
  Atenção: A prova é baseada no livro, não nas apresentações
\end{frame}

%------------------------------------------------------------------------------%
\end{document}
%------------------------------------------------------------------------------%
